%! AP Statistics — Unit 1, Lesson 1.6 — Warm-Up
\documentclass[10pt]{article}
\usepackage{apstats-article}
\usepackage{apstats-boxes}

%=============================================================================
\begin{document}

\pageheader{Unit 1, Lesson 1.6}{Warm-Up}
\namedateperiod

\vspace{0.18in}

\noindent In Lesson 1.5 you constructed histograms and dotplots for quantitative
data. Today you will practice \textit{describing} those displays --- turning
what you see in a graph into precise, complete written analysis.

\vspace{0.12in}

\noindent The histogram below (already drawn for you) shows the \textbf{number
of hours per week} that 30 students spend on homework.

\vspace{0.1in}

\begin{center}
\small
\renewcommand{\arraystretch}{1.5}
\begin{tabular}{@{} >{\bfseries}p{0.22\linewidth} p{0.70\linewidth} @{}}
\rowcolor{sky}
\textbf{Hours / Week} & \textbf{Frequency (number of students)} \\
\midrule
0--under 4  & \textcolor{navy}{\rule{1.5cm}{10pt}} \quad 3 \\
4--under 8  & \textcolor{navy}{\rule{4.5cm}{10pt}} \quad 9 \\
8--under 12 & \textcolor{navy}{\rule{5.5cm}{10pt}} \quad 11 \\
12--under 16 & \textcolor{navy}{\rule{3.0cm}{10pt}} \quad 6 \\
16--under 20 & \textcolor{navy}{\rule{0.5cm}{10pt}} \quad 1 \\
\end{tabular}
\end{center}

\vspace{0.18in}

\begin{enumerate}

  \item Looking at the bar lengths above, describe the overall \textit{shape}
        of this distribution using one of these words: symmetric, skewed right,
        skewed left, bimodal, uniform. Explain your choice.\\[6pt]
        \writeline\\[1pt]\writeline

  \item Estimate a \textit{typical} (center) value for homework hours per week
        for these 30 students. What in the display supports your estimate?\\[6pt]
        \writeline\\[1pt]\writeline

  \item Does the distribution appear to have any \textit{outliers} --- values
        that fall far from the overall pattern? Which interval(s)?\\[6pt]
        \writeline\\[1pt]\writeline

\end{enumerate}

\vspace{0.1in}

\begin{tcolorbox}[colback=greenbg, colframe=greenacc, boxrule=0.8pt, arc=2mm,
    left=4mm, right=4mm, top=3mm, bottom=3mm]
\small
\begin{itemize}[leftmargin=1.3em, itemsep=5pt]
  \item Today we formalize a four-part framework --- \textbf{SOCS} --- for
        describing any quantitative distribution completely.
  \item SOCS stands for \textbf{S}hape, \textbf{O}utliers, \textbf{C}enter,
        and \textbf{S}pread. The AP exam requires all four components.
  \item Every sentence in your description must include \textit{context}
        (the variable name and population).
\end{itemize}
\end{tcolorbox}

\vspace{0.1in}

\noindent\textcolor{slate}{\small One question you have going into today's lesson,
or something you're curious about:}\\[8pt]
\writeline

\end{document}
